\chapter{Cadrage de besoin et analyse}

\section{Introduction}
Ce chapitre a pour objectif de présenter le cadrage complet du projet 
développé dans le cadre du stage chez Bewize. Il commence par une analyse 
approfondie du contexte et des besoins, puis décrit l'étude de l'existant 
et enfin la solution proposée pour répondre aux attentes de Bewize et 
de ses utilisateurs. Ce cadrage constitue la base indispensable pour 
aborder sereinement les phases de conception et de développement 
présentées dans les chapitres suivants.

%----------------------------------------------------------
\section{Contexte du projet}
%----------------------------------------------------------

\subsection{Présentation du besoin métier}
Bewize est une entreprise proposant une application mobile destinée à 
ses utilisateurs. Dans le cadre de son développement, Bewize souhaite 
améliorer le suivi des retours de ses utilisateurs. Aujourd'hui, aucun 
mécanisme simple et intégré ne permet aux utilisateurs de transmettre 
directement leurs feedbacks ou réclamations depuis l'application mobile.

Cette absence de canal de communication directe entre les utilisateurs 
et l'équipe Bewize représente un manque important. Les utilisateurs 
insatisfaits n'ont pas de moyen simple pour signaler un problème, et 
l'équipe interne ne dispose pas d'une vue centralisée des retours 
utilisateurs. Cela peut entraîner une dégradation progressive de la 
satisfaction client et une perte d'informations précieuses pour 
l'amélioration du produit.

\subsection{Enjeux du projet}
Les enjeux de ce projet sont multiples et touchent à la fois à la qualité 
du produit, à la satisfaction client et à l'efficacité opérationnelle 
de l'équipe Bewize :

\begin{itemize}
    \item \textbf{Enjeu produit :} améliorer la qualité de l'application 
    mobile en collectant les retours directs des utilisateurs et en les 
    intégrant dans le processus d'amélioration continue.
    
    \item \textbf{Enjeu satisfaction client :} offrir aux utilisateurs 
    un canal simple, rapide et accessible pour exprimer leurs feedbacks 
    ou réclamations, renforçant ainsi leur sentiment d'être écoutés.
    
    \item \textbf{Enjeu opérationnel :} centraliser les retours 
    utilisateurs dans le dashboard monitor interne afin de faciliter 
    leur traitement par l'équipe Bewize.
    
    \item \textbf{Enjeu technique :} intégrer cette nouvelle fonctionnalité 
    de manière cohérente avec l'architecture existante de Bewize.
\end{itemize}

\subsection{Périmètre du projet}
Le projet est encadré par le Product Manager de Bewize, en collaboration 
avec le Tech Lead et l'équipe de développement. Il s'inscrit dans un 
stage de deux mois et couvre les phases suivantes :

\begin{enumerate}
    \item Immersion agile et formation Scrum
    \item Cadrage, conception et développement de la fonctionnalité
    \item Tests et validation du processus complet
    \item Découverte du marketing digital chez Bewize
\end{enumerate}

%----------------------------------------------------------
\section{Analyse des besoins}
%----------------------------------------------------------

\subsection{Besoins fonctionnels}
Les besoins fonctionnels décrivent l'ensemble des fonctionnalités que 
le système doit offrir pour répondre aux attentes des utilisateurs et 
de l'équipe Bewize. Après une analyse approfondie du contexte et des 
attentes métier, les besoins fonctionnels suivants ont été identifiés :

\begin{table}[H]
\centering
\caption{Tableau des besoins fonctionnels}
\label{tab:besoins_fonctionnels}
\renewcommand{\arraystretch}{1.8}
\begin{tabular}{|c|p{4cm}|p{7cm}|}
\hline
\textbf{ID} & \textbf{Besoin} & \textbf{Description} \\
\hline
BF01 & Accès à la fonctionnalité & L'utilisateur doit pouvoir accéder 
facilement à la fonctionnalité de feedback depuis l'application mobile 
via un bouton dédié. \\
\hline
BF02 & Redirection WebView & L'utilisateur est redirigé vers une page 
WebView simple et claire après avoir cliqué sur le bouton. \\
\hline
BF03 & Formulaire de saisie & L'utilisateur peut saisir son message de 
feedback ou de réclamation dans un formulaire simple et intuitif. \\
\hline
BF04 & Envoi des données & Les données saisies par l'utilisateur sont 
envoyées via API au back-end Bewize. \\
\hline
BF05 & Enregistrement & Les informations sont enregistrées dans la base 
de données Bewize pour consultation ultérieure. \\
\hline
BF06 & Consultation dashboard & L'équipe interne Bewize peut consulter 
les feedbacks et réclamations reçus depuis le dashboard monitor. \\
\hline
BF07 & Centralisation & Toutes les données utilisateurs et réclamations 
sont centralisées dans un seul espace accessible à l'équipe. \\
\hline
\end{tabular}
\end{table}

Le tableau \ref{tab:besoins_fonctionnels} présente les sept besoins 
fonctionnels identifiés lors de la phase d'analyse. Ces besoins couvrent 
l'ensemble du parcours utilisateur, depuis l'accès à la fonctionnalité 
depuis l'application mobile jusqu'à la consultation des retours par 
l'équipe interne Bewize dans le dashboard monitor.

\subsection{Besoins non fonctionnels}
Les besoins non fonctionnels définissent les contraintes de qualité, 
de performance et de sécurité que le système doit respecter pour 
garantir une expérience utilisateur optimale et une solution fiable :

\begin{table}[H]
\centering
\caption{Tableau des besoins non fonctionnels}
\label{tab:besoins_non_fonctionnels}
\renewcommand{\arraystretch}{1.8}
\begin{tabular}{|c|p{4cm}|p{7cm}|}
\hline
\textbf{ID} & \textbf{Besoin} & \textbf{Description} \\
\hline
BNF01 & Performance & La page WebView doit se charger rapidement pour 
ne pas nuire à l'expérience utilisateur. \\
\hline
BNF02 & Simplicité & L'interface doit être simple, claire et intuitive 
pour tout type d'utilisateur. \\
\hline
BNF03 & Fiabilité & Les données envoyées par l'utilisateur doivent être 
transmises de manière fiable sans perte d'information. \\
\hline
BNF04 & Sécurité & Les données personnelles des utilisateurs doivent 
être protégées et traitées de manière sécurisée. \\
\hline
BNF05 & Compatibilité & La solution doit être compatible avec 
l'application mobile Bewize existante et son architecture technique. \\
\hline
BNF06 & Maintenabilité & Le code développé doit être propre, documenté 
et facile à maintenir par l'équipe de développement Bewize. \\
\hline
\end{tabular}
\end{table}

Le tableau \ref{tab:besoins_non_fonctionnels} récapitule les six besoins 
non fonctionnels retenus pour ce projet. Ces contraintes de qualité 
sont essentielles pour garantir que la solution développée soit non 
seulement fonctionnelle, mais également performante, sécurisée et 
maintenable sur le long terme.

%----------------------------------------------------------
\section{Étude de l'existant}
%----------------------------------------------------------

\subsection{Présentation de l'existant}
Avant de concevoir la solution, il est essentiel d'analyser l'existant 
afin de comprendre le contexte technique dans lequel la fonctionnalité 
va s'intégrer. Bewize dispose déjà d'un écosystème technique composé 
de plusieurs éléments interdépendants :

\begin{itemize}
    \item \textbf{Application mobile Bewize :} l'application mobile est 
    déjà développée et utilisée par les utilisateurs finaux. Elle 
    constitue le point d'entrée principal de la nouvelle fonctionnalité.
    
    \item \textbf{Dashboard monitor interne :} interface web développée 
    en interne par Bewize, permettant à l'équipe de gérer et consulter 
    les données liées aux utilisateurs. Ce dashboard sera enrichi pour 
    afficher les feedbacks et réclamations collectés.
    
    \item \textbf{APIs Spring Boot :} Bewize utilise déjà des APIs 
    développées avec Spring Boot pour la communication entre les 
    différentes parties de son système.
    
    \item \textbf{Base de données Bewize :} une base de données 
    centralisée existe déjà pour stocker les données utilisateurs 
    et applicatives.
\end{itemize}

\subsection{Limites de l'existant}
Malgré cet écosystème technique déjà en place, plusieurs limites ont 
été identifiées lors de la phase d'analyse :

\begin{itemize}
    \item Absence de canal direct de communication entre les utilisateurs 
    et l'équipe Bewize depuis l'application mobile.
    \item Aucun mécanisme de collecte structurée des feedbacks et 
    réclamations n'est intégré à l'application mobile.
    \item L'équipe interne ne dispose pas d'une vue centralisée et 
    organisée des retours utilisateurs.
    \item Le traitement des réclamations est actuellement manuel et 
    non systématisé.
\end{itemize}

Ces limites justifient pleinement le développement de la nouvelle 
fonctionnalité et démontrent la valeur ajoutée qu'elle apportera 
à l'écosystème Bewize.

%----------------------------------------------------------
\section{Solution proposée}
%----------------------------------------------------------

\subsection{Description de la solution}
La solution proposée consiste à développer une fonctionnalité complète 
de recueil de feedbacks et de réclamations, intégrée directement dans 
l'application mobile Bewize. Cette solution repose sur l'utilisation 
d'une page WebView accessible depuis l'application mobile, communiquant 
avec le back-end via des APIs Spring Boot, et dont les données sont 
visualisées dans le dashboard monitor.

\subsection{Parcours utilisateur}
Le parcours utilisateur a été conçu pour être le plus simple et fluide 
possible. Il se déroule en sept étapes :

\begin{enumerate}
    \item L'utilisateur ouvre l'application mobile Bewize et navigue 
    jusqu'au bouton de feedback ou réclamation.
    \item Il clique sur ce bouton dédié, clairement identifiable 
    dans l'interface.
    \item Une page WebView s'ouvre automatiquement, affichant un 
    formulaire simple et clair.
    \item L'utilisateur remplit le formulaire en saisissant son message 
    ainsi que les informations nécessaires.
    \item Il soumet le formulaire en cliquant sur le bouton d'envoi.
    \item Les données sont transmises via API au back-end Spring Boot, 
    qui les traite et les enregistre dans la base de données Bewize.
    \item L'équipe interne Bewize consulte les feedbacks et réclamations 
    depuis le dashboard monitor et procède à leur traitement.
\end{enumerate}

\subsection{Fonctionnalités détaillées}
La fonctionnalité développée dans le cadre de ce stage couvre les 
aspects suivants, organisés par composant :

\begin{itemize}
    \item \textbf{Côté application mobile :}
    \begin{itemize}
        \item Ajout d'un bouton d'accès au feedback et à la réclamation
        \item Intégration d'une WebView pour afficher le formulaire
        \item Gestion de la navigation entre l'application et la WebView
    \end{itemize}
    
    \item \textbf{Côté page WebView :}
    \begin{itemize}
        \item Affichage d'un formulaire simple et responsive
        \item Saisie du type de retour (feedback ou réclamation)
        \item Validation et envoi des données via API
    \end{itemize}
    
    \item \textbf{Côté back-end Spring Boot :}
    \begin{itemize}
        \item Création des endpoints API nécessaires
        \item Réception et traitement des données envoyées
        \item Enregistrement dans la base de données Bewize
    \end{itemize}
    
    \item \textbf{Côté dashboard monitor :}
    \begin{itemize}
        \item Affichage de la liste des feedbacks et réclamations reçus
        \item Consultation du détail de chaque retour
        \item Possibilité de traitement et suivi par l'équipe interne
    \end{itemize}
\end{itemize}

%----------------------------------------------------------
\section{Planning du stage}
%----------------------------------------------------------

Le stage de deux mois chez Bewize s'étend sur huit semaines. Le planning 
a été organisé de manière à couvrir toutes les phases du projet, de 
l'immersion initiale jusqu'à la restitution finale.

\begin{table}[H]
\centering
\caption{Planning détaillé du stage}
\label{tab:planning}
\renewcommand{\arraystretch}{1.8}
\begin{tabular}{|c|p{3.5cm}|p{5cm}|p{3cm}|}
\hline
\textbf{Semaine} & \textbf{Phase} & \textbf{Objectifs principaux} & 
\textbf{Encadrement} \\
\hline
Semaine 1 & Immersion agile & Découverte de Bewize, formation à 
l'agilité, introduction au framework Scrum & 
Product Manager et équipe Bewize \\
\hline
Semaine 2 & Cadrage produit & Analyse du besoin business, étude de 
l'existant, conception fonctionnelle & Product Manager \\
\hline
Semaine 3 & Conception technique & Définition de l'architecture, 
début du développement front-end et back-end & Tech Lead et équipe dev \\
\hline
Semaine 4 & Développement & Construction des composants, premières 
intégrations & Tech Lead et équipe dev \\
\hline
Semaine 5 & Intégration & Connexion base de données, liaisons API & 
Tech Lead et équipe dev \\
\hline
Semaine 6 & Tests et validation & Tests fonctionnels, correction des 
anomalies & Product Manager et Tech Lead \\
\hline
Semaine 7 & Marketing digital & Découverte des pratiques marketing 
chez Bewize & Product Manager \\
\hline
Semaine 8 & Clôture & Synthèse, documentation, restitution finale & 
Product Manager et équipe Bewize \\
\hline
\end{tabular}
\end{table}

Le tableau \ref{tab:planning} présente la répartition des activités 
sur les huit semaines du stage. Ce planning équilibré permet de 
consacrer les premières semaines à la compréhension du contexte et 
à la conception, les semaines centrales au développement et à 
l'intégration, et les dernières semaines à la validation, aux tests 
et à la découverte du marketing digital chez Bewize.

%----------------------------------------------------------
\section{Conclusion}
%----------------------------------------------------------

Ce chapitre a permis de définir le cadre complet du projet, en analysant 
le contexte, les besoins fonctionnels et non fonctionnels, l'existant 
technique de Bewize et la solution proposée. L'analyse des besoins a 
mis en évidence sept besoins fonctionnels et six besoins non fonctionnels 
qui guideront l'ensemble du développement. L'étude de l'existant a 
confirmé la pertinence et la valeur ajoutée de la solution à développer. 
Cette analyse constitue la base indispensable pour aborder sereinement 
les phases de conception et de développement présentées dans les 
chapitres suivants.
